\documentclass[11pt]{article}
\usepackage[a4paper,margin=2.2cm,top=2.6cm]{geometry}
\usepackage{amsmath,amssymb}
\usepackage{booktabs}
\usepackage{fancyhdr}
\usepackage{enumitem}
\usepackage[hidelinks]{hyperref}

% Edit these two macros to match the program's header.
\newcommand{\schoolname}{Efi Arazi School of Computer Science \quad Machine Learning and Data Science M.Sc. Program}
\newcommand{\proposalname}{MLDS Project Proposal}

\pagestyle{fancy}
\fancyhf{}
\renewcommand{\headrulewidth}{0.4pt}
\fancyhead[L]{\footnotesize\schoolname}
\fancyfoot[R]{\footnotesize\proposalname\quad Page \thepage}

\setlength{\parindent}{0pt}
\setlength{\parskip}{0.55em}
\setlist{itemsep=0.2em,topsep=0.3em,leftmargin=1.4em}
\newcommand{\RQ}[1]{\par\medskip\noindent\textbf{Research question:} #1\par\medskip}
\newcommand{\E}{\mathbb{E}}

\begin{document}

{\Large\bfseries \proposalname\par}
\vspace{0.2em}
{\large\bfseries When Does a Label-Free Verifier Ensemble Pick the Wrong Signal?\\
Correlated Weak Verifiers and Shared Nuisance Factors in Best-of-$N$ Selection\par}
\vspace{0.6em}

\textbf{Date:} November 2026\\
\textbf{Mentors:} TBD

\section*{Background}

Best-of-$N$ selection asks a language model for $N$ candidate answers and lets a set of
imperfect \emph{verifiers} (LLM judges, reward models, process reward models) choose
among them. Combining many weak verifiers has been shown to close much of the gap
between a model's first answer and its best-of-$N$ ceiling~\cite{weaver}, but the
combination step normally needs labelled questions to learn verifier weights. FUSE
\cite{fuse} removes the labels: it treats each binarised verifier as a noisy classifier,
estimates every verifier's sensitivity and specificity by the method of moments of
Jaffe et al.~\cite{jaffe}, and ranks responses by the resulting posterior. The practical
appeal is large, since no one has to label anything, and so is the risk: a label-free
method has no signal telling it when its own assumptions have broken.

\textbf{Preliminaries.} For a question, let $Y_i\in\{0,1\}$ be the unobserved correctness
of response $i=1,\dots,N$ and $V_{ij}\in\{0,1\}$ the binarised output of verifier
$j=1,\dots,m$. The moment estimator assumes that verifiers are conditionally independent
given $Y$. Then the off-diagonal covariance of $V$ is exactly rank one,
$\Sigma_{jk}=\lambda_j\lambda_k$ for $j\neq k$, where $\lambda_j\propto
(\psi_j+\eta_j-1)$ is the informativeness of verifier $j$ (its sensitivity $\psi_j$ plus
specificity $\eta_j$ minus one), and $\lambda$ is identifiable up to
sign from any three verifiers. The assumption can fail in the way that matters. Suppose
a nuisance factor $U_i$ (response length, formatting, a shared backbone model, a shared
training set) drives a block $B$ of verifiers and is unrelated to $Y$. Then
$\Sigma_{jk}$ is large for $j,k\in B$, the block's rank-one fit is excellent, and a
label-free method cannot distinguish the factor $U$ from the factor $Y$: the only thing
it sees is a strong shared signal. Which factor it labels ``correct'' depends on how large
the block is, not on which factor is true.

\textbf{Pilot evidence.} A from-scratch reimplementation of FUSE and its Best-of-$N$
experiments on the authors' released data (12 settings) reproduces the paper to within
1.5--2 points on most of them. On MATH500 with 8B-generated candidates it does not: the
best configuration reaches 66.0\% against the paper's 83.9\%. Nine of the 16 reward models
are near-random (selection accuracy about 54\% against a 50\% single-sample baseline) but
mutually correlated at $|\rho|\approx 0.5$, while the two strongest verifiers are weakly
correlated with everything. The rank-one fit locks onto the correlated block and, estimated
on all questions at once, ranks it best and the two strongest verifiers worst. Averaging
correlated blocks into one meta-verifier recovers MATH500-8B to 79.3\% but costs
GPQA-70B 2--10 points, because the correlated LM judges there are informative. Correlation
among verifiers is therefore ambiguous: shared signal and shared noise look the same.

\begin{center}\small
\begin{tabular}{lrrr}\toprule
Setting (query mode) & FUSE (ours) & + clustering & Paper \\\midrule
MATH500-8B & 66.0 & 79.3 & 83.9 \\
GPQA-70B & 60.5 & 58.0 & 66.8 \\
MMLU-8B & 80.5 & 82.6 & 84.9 \\\bottomrule
\end{tabular}\\[0.2em]
\textit{Table 1. Pilot results from the reimplementation: one correlated-block failure and its cost.}
\end{center}

\textbf{Core: when it fails, whether we can tell, and what repairs it.} The core has
five movements, in three tiers: (a)--(c) are mandatory, (d) and (e) secondary,
and (f) exploratory.
\begin{enumerate}[label=\alph*)]
\item \textbf{Theory.} Model the verifier covariance as a sum of two low-rank terms,
$\Sigma_{\text{off}}=\lambda\lambda^\top+\mu\mu^\top$ ($\mu$ the loading on $U$, supported on
$B$). Derive when the best rank-one fit aligns with $\mu$ rather than $\lambda$ as a
function of block size $|B|$, within-block loading, and the informative verifiers'
strength, and state the identifiability limit: a second factor of equal size cannot be
resolved from second moments alone.
\item \textbf{Controlled simulation with known ground truth.} Generate verifier matrices
from the two-factor model with planted $B$, sweeping $|B|$ and $\mu$, and compare FUSE,
Jaffe et al., Dawid--Skene, the naive ensemble and majority vote on selection accuracy,
reported with a matched arm in which $\mu=0$ so that no remedy is rewarded merely for
being conservative.
\item \textbf{A label-free failure detector.} A statistic that fires when the data look
two-factor rather than one-factor: the spectral gap of the verifier correlation matrix
after removing the leading component, and a triplet conditional-independence test. It is
evaluated as a detector, by its false-alarm rate on settings where FUSE is fine and its
detection rate on settings where it fails, in simulation and on real data. The question
is not whether it explains the failure but whether it would have warned the user.
\item \textbf{Nuisance audit on real data.} On the 12 released settings, measure how much
of the inter-verifier correlation is explained by observable covariates (response length,
presence of a final answer, answer format, generator identity) and whether partialling
them out restores the one-factor structure.
\item \textbf{Remedy comparison under one protocol.} Verifier clustering, nuisance
residualisation, an explicit rank-two model, and abstention to the naive ensemble when
the detector fires; each judged on both the failing settings and the settings where
correlated verifiers are informative.
\item \textbf{Exploratory: identifying the correct factor.} Whether a handful of labelled
questions (1--5\%) suffice to decide which of two factors is correctness. Reported as an
open problem, not a result.
\end{enumerate}

\textbf{Evaluation discipline.} The aim is a statement about a method, so the ways to
overstate it are addressed directly.
\emph{Tuning on the test settings.} Remedy hyperparameters (clustering threshold, detector
cutoff) are fixed on the simulation arm before any real-data run, and every real setting
is reported, not the best cell per benchmark.
\emph{Label leakage into a label-free method.} The detector and every remedy are computed
without labels; labels enter only in scoring, and splits are by question.
\emph{Selection among many configurations.} With two modes, two predictors and several
thresholds, differences of 1--2 points are within noise, so selection accuracy is reported
with paired bootstrap intervals over questions and compared with the naive ensemble as an
explicit floor.
\emph{Confounding.} The nuisance audit is observational; a covariate that explains
correlation is evidence, not proof, of a mechanism.
\emph{Size and power.} A remedy that repairs MATH500-8B but harms GPQA-70B is judged
against both, using accuracy change per setting rather than the mean alone.

\section*{Objectives \& Deliverables}

\textbf{Objectives}
\begin{itemize}
\item Formalise the two-factor failure of moment-based label-free ensembling and state
when it is, and is not, identifiable.
\item Build a simulation framework with planted nuisance blocks and known ground truth,
including a matched null arm.
\item Design and validate a label-free detector of the failure, and quantify its false-alarm
and detection rates. This is the project's central contribution.
\item Compare the available remedies head-to-head on size and power, on simulated and
released data.
\item Audit the 12 released settings for the nuisance factors behind the failure.
\end{itemize}

\textbf{Deliverables}
\begin{itemize}
\item A tested Python library extending the existing FUSE reimplementation with the
detector, the remedies, and the simulator.
\item A benchmark table over all 12 settings and the simulation grid, with paired
bootstrap intervals.
\item An instructional notebook that reproduces the MATH500-8B failure in one sitting.
\item A final report stating, falsifiably, when label-free ensembling can be trusted, which
remedy is safe to ship, and where the problem remains open.
\end{itemize}

\section*{Requirements}

Probability and statistics at the level of an advanced statistics course; Python; some
familiarity with language-model evaluation. Expected workload is 1--2 days a week for an
academic year.

\section*{Resources}

Data: the verifier score matrices released with FUSE (HLE, IMO Shortlist) and the Weaver
datasets, used as they are; no model needs to be run. Everything runs on a laptop. The
starting code base (FUSE reimplementation, baselines, results for all 12 settings) is
available to the student. Academic project; any developed IP will be addressed when such
occurs.

\begin{thebibliography}{9}\small
\bibitem{fuse} Lee, Ma, Zhao, et al. Fully Unsupervised Score Ensembling (FUSE). arXiv:2604.18547.
\bibitem{weaver} Saad-Falcon, J., et al. (2025). Shrinking the Generation-Verification Gap with Weak Verifiers (Weaver).
\bibitem{jaffe} Jaffe, A., Nadler, B., Kluger, Y. (2015). Estimating the accuracies of multiple classifiers without labeled data. AISTATS.
\bibitem{ds} Dawid, A. P., Skene, A. M. (1979). Maximum likelihood estimation of observer error-rates using the EM algorithm. J. Royal Statistical Society C 28(1), 20--28.
\bibitem{parisi} Parisi, F., Strino, F., Nadler, B., Kluger, Y. (2014). Ranking and combining multiple predictors without labeled data. PNAS 111(4), 1253--1258.
\bibitem{lightman} Lightman, H., et al. (2023). Let's Verify Step by Step. arXiv:2305.20050.
\end{thebibliography}

\end{document}
